\chapter{Appendix}

\section{Complete Module Listings}

This appendix contains complete code listings for all processor modules.

\subsection{Top-Level Module: single\_cycle.sv}

Due to length, the complete code is available in the source directory. Key sections include:

\begin{itemize}
    \item Module declaration with all I/O ports
    \item Internal signal declarations
    \item PC update logic
    \item Instruction fetch
    \item Immediate generation
    \item ALU operand selection
    \item Branch and jump logic
    \item Memory-mapped I/O handling
    \item Module instantiations
\end{itemize}

\subsection{Register File Module}

The register file implementation is shown in Chapter 4. Key features:

\begin{itemize}
    \item 32 registers of 32 bits each
    \item Asynchronous read, synchronous write
    \item Register x0 hardwired to zero
\end{itemize}

\subsection{ALU Module}

The ALU supports 10 operations as detailed in Chapter 4:

\begin{itemize}
    \item ADD, SUB, AND, OR, XOR
    \item SLL, SRL, SRA
    \item SLT, SLTU
\end{itemize}

\subsection{Control Unit Module}

The control unit decodes instructions and generates control signals for:

\begin{itemize}
    \item R-type instructions (7 opcodes)
    \item I-type instructions (multiple variants)
    \item S-type, B-type, U-type, J-type instructions
\end{itemize}

\section{File Structure}

\subsection{Source Files (00\_src/)}

\begin{itemize}
    \item \texttt{single\_cycle.sv}: Top-level module
    \item \texttt{register\_file.sv}: Register file
    \item \texttt{alu.sv}: Arithmetic Logic Unit
    \item \texttt{control\_unit.sv}: Control unit
    \item \texttt{imem.sv}: Instruction memory
    \item \texttt{dmem.sv}: Data memory
\end{itemize}

\subsection{Testbench Files (01\_bench/)}

\begin{itemize}
    \item \texttt{tbench.sv}: Top-level testbench
    \item \texttt{driver.sv}: Stimulus driver
    \item \texttt{scoreboard.sv}: Result checker
    \item \texttt{tlib.svh}: Utility library
\end{itemize}

\subsection{Simulation Directories}

\begin{itemize}
    \item \texttt{10\_sim/}: Verilator simulation files
    \item \texttt{11\_xm/}: Xcelium simulation files
\end{itemize}

\section{Memory Map Details}

\subsection{Complete Address Map}

Table~\ref{tab:complete_memory_map} shows the complete memory map.

\begin{table}[h]
\centering
\caption{Complete Memory Map}
\label{tab:complete_memory_map}
\begin{tabular}{lll}
\toprule
\textbf{Base} & \textbf{Top} & \textbf{Description} \\
\midrule
0x1001\_1000 & 0xFFFF\_FFFF & Reserved \\
0x1001\_0000 & 0x1001\_0FFF & Switches (read-only) \\
0x1000\_5000 & 0x1000\_FFFF & Reserved \\
0x1000\_4000 & 0x1000\_4FFF & LCD Control \\
0x1000\_3000 & 0x1000\_3FFF & HEX4-7 Displays \\
0x1000\_2000 & 0x1000\_2FFF & HEX0-3 Displays \\
0x1000\_1000 & 0x1000\_1FFF & Green LEDs \\
0x1000\_0000 & 0x1000\_0FFF & Red LEDs \\
0x0000\_0800 & 0x0FFF\_FFFF & Reserved \\
0x0000\_0000 & 0x0000\_07FF & Data Memory (2KB) \\
\bottomrule
\end{tabular}
\end{table}

\section{Instruction Encoding}

\subsection{RISC-V Instruction Formats}

All RISC-V instructions are 32 bits wide with the following formats:

\begin{itemize}
    \item \textbf{R-type}: opcode, rd, funct3, rs1, rs2, funct7
    \item \textbf{I-type}: opcode, rd, funct3, rs1, immediate[11:0]
    \item \textbf{S-type}: opcode, imm[4:0], funct3, rs1, rs2, imm[11:5]
    \item \textbf{B-type}: opcode, imm[11], imm[4:1], funct3, rs1, rs2, imm[10:5], imm[12]
    \item \textbf{U-type}: opcode, rd, imm[31:12]
    \item \textbf{J-type}: opcode, rd, imm[19:12], imm[11], imm[10:1], imm[20], imm[30:21], imm[31]
\end{itemize}

\section{Control Signal Truth Table}

Due to space limitations, the complete truth table is available in the source code comments. Key control signal combinations:

\begin{itemize}
    \item R-type: reg\_write=1, alu\_src\_a=00, alu\_src\_b=00
    \item I-type: reg\_write=1, alu\_src\_a=00, alu\_src\_b=01
    \item Load: reg\_write=1, mem\_read=1, mem\_to\_reg=01
    \item Store: mem\_write=1, alu\_src\_a=00, alu\_src\_b=01
    \item Branch: branch=1, alu\_src\_a=00, alu\_src\_b=00
    \item Jump: jump=1, reg\_write=1, mem\_to\_reg=10
\end{itemize}

\section{Simulation Commands}

\subsection{Verilator}

\begin{verbatim}
cd 10_sim
verilator --cc --exe --build --top-module tbench -f ./flist
./obj_dir/Vtbench
\end{verbatim}

\subsection{Xcelium}

\begin{verbatim}
cd 11_xm
module load xcelium
xrun -64bit -access +rwc -f ./flist
simvision wave.shm
\end{verbatim}

\section{Testing Checklist}

Before submission, verify:

\begin{enumerate}
    \item All files compile without errors
    \item All modules instantiated correctly
    \item Memory addresses match specification
    \item I/O mapping correct
    \item Test suite passes
    \item File structure matches requirements
    \item Documentation complete
\end{enumerate}

\section{References}

All source code and testbenches follow:

\begin{itemize}
    \item RISC-V Instruction Set Manual, Volume I: User-Level ISA
    \item SystemVerilog IEEE 1800-2017 standard
    \item Milestone 2 specification document
    \item TA-provided testbench requirements
\end{itemize}

